\chapter{\MakeUppercase{Project implementation}}
\justifying
\section{Backend Implementation}
In the server/app.py file, we have the entry point for our backend system. Here, we set up our FastAPI application and enable cross-origin settings, load model databases (also known as checkpoints), set up the feature converter, provide access to the database, and create the set of services required at runtime (enrollment, authentication, and handling of challenges).

By using a startup model approach for our API, handler functions only execute business logic, meaning that they will not have to re-create any models. These are built using Pydantic request and response objects, which is important for ensuring that the application enforces field requirements, types/constraints on fields, and shape validations of the returned payload (e.g. samples sent by clients).

Furthermore, Pydantic request and response definitions will lead to improved quality of OpenAPI documentation generated by our backend application and will be available via the /docs endpoint.

\section{Enrollment Module}
Auth/enrollment.py has the enrollment implementation, where it creates device IDs, gets/normalizes raw samples, transforms them into features, generates source embeddings, builds template banks, saves template bundles, and records metadata. Metadata contains the following:
\begin{itemize}
	\item Device name,
	\item Number of samples,
	\item Active sources,
	\item Feature dimension,
	\item Source profiles,
	\item Source weights,
	\item Template strategy and template counts,
	\item Drift model parameters,
	\item Attacker-model notes.
\end{itemize}

This is one of the stronger parts in the overall implementation because having metadata makes each enrollment interpretable and easily reviewable.

\section{Authentication Module}
The file auth/authentication.py is where the implementation of Authentication is housed. In this file, the source of authentication is validated, runtime profiles are computed, matches are checked against possible mismatches, runtime embeddings are created, scores are generated against stored template banks and confidence-band plus margin logic is applied. The module allows for rolling re-enrollment after strong matches have been identified in the case of drift updates being enabled.

One special note in the development of the module was to use a top-k scoring methodology when comparing the scores in the template banks. Rather than using just a single best or a single average score from the matching templates, the module creates an average score of the top x number of template similarities. This provides less influence on the score of one outlier template.

\section{Challenge-Response Module}
The challenge-response logic is implemented in \texttt{auth/challenge\_response.py}. A challenge contains a unique challenge ID, device ID, random nonce, creation time, and expiry. The stable feature bytes are derived by coarse quantization of the embedding, then combined with the nonce and server secret through \ac{hkdf} to derive the MAC key. The final response is an \ac{hmac} over challenge data. This prevents straightforward replay because each challenge is single-use and nonce-bound.

\section{Database Implementation}
The persistence layer is implemented with SQLAlchemy in \texttt{db/models.py}. There are three tables:
\begin{itemize}
	\item \texttt{Device}: stores device ID, name, embedding path, metadata, and timestamp,
	\item \texttt{Challenge}: stores active challenge state and expiry,
	\item \texttt{AuditLog}: stores event details for important actions.
\end{itemize}

The database is SQLite by default, which is reasonable for a capstone prototype and simplifies setup. The code also keeps the schema narrow, which reduces maintenance complexity.

\section{Frontend Implementation}
The frontend application under \texttt{frontend/src/} provides a simple but effective control surface. The home page frames the project as ``noise-based device verification from camera and microphone fingerprints.'' Dedicated pages are available for enrollment, authentication, and device viewing. The frontend service layer centralizes HTTP calls to the backend and defines typed request and response interfaces.

\section{Configuration and Operational Controls}
Configuration values in \texttt{config.py} make the implementation flexible without being over-engineered. Important controls include:
\begin{itemize}
	\item model checkpoint paths,
	\item confidence thresholds,
	\item identification margin,
	\item source weights,
	\item drift-update parameters,
	\item challenge expiry,
	\item demo-mode source limits,
	\item optional API key and rate limiting settings.
\end{itemize}

Such configurability is valuable in a thesis setting because it makes experimentation explicit and reproducible.

\section{Illustrative Code Snippet}
The following excerpt illustrates the style of the implementation by showing a representative configuration block:

\begin{lstlisting}[language=Python]
AUTH_CONFIDENCE_STRONG = 0.97
AUTH_CONFIDENCE_UNCERTAIN = 0.92
AUTH_IDENTIFICATION_MARGIN = 0.02
AUTH_TEMPLATE_CHUNK_SIZE = 5
AUTH_MAX_TEMPLATES_PER_SOURCE = 4
AUTH_DRIFT_EMA_ALPHA = 0.2
AUTH_SOURCE_WEIGHTS = {
    "camera": 0.7,
    "microphone": 0.3,
}
\end{lstlisting}

This reflects the overall implementation philosophy of the project: keep the runtime rules explicit and inspectable rather than burying key decisions inside opaque code paths.
